%%==========================================================================%%
%% methods.tex -- "Methods". Subsection structure and boilerplate modeled on %%
%% the reference Scientific Data paper; prose paraphrased/genericized for     %%
%% whole-body pediatric MRI. REWRITE and verify before submission.           %%
%% [TBD: ...] = fill in our specifics.                                        %%
%%==========================================================================%%

\section*{Methods}\label{sec:methods}

\subsection*{Ethics statement}\label{subsec:ethics}
% Reusable skeleton: approving body + accordance + legal basis for sharing +
% consent / de-identification status. ADAPT to our institution (e.g., Stanford
% IRB + HIPAA) -- the reference uses Danish/EU GDPR wording, which does NOT
% apply to us. See the de-identification TODO in CLAUDE.md.
The study was approved by [TBD: named institutional review board / ethics
committee, approval number], and all procedures were carried out in accordance
with the relevant guidelines and regulations [TBD: e.g., the Declaration of
Helsinki]. The data are shared under [TBD: legal basis / consent model], and
[TBD: informed consent / waiver and consent for public sharing] was obtained.
[TBD: state the de-identification status of the released data and any resulting
access conditions --- see Usage Notes.]

\subsection*{Participants}\label{subsec:cohort}
% Reusable template: N, sex split, age (mean +/- SD), recruitment source,
% acquisition date range, exclusions with counts.
The released cohort comprises 250 whole-body PET/MRI examinations from 134
patients aged 1 to 25 years at the time of imaging (median 14 years; sex split
[TBD: $N$ female, $N$ male]).
Subjects were retrospectively identified at [TBD: site(s)], with examinations
acquired between [TBD: start] and [TBD: end].
The released examinations form the union of three annotated subsets, which
partially overlap (a single examination may belong to more than one subset, so
the subset counts do not sum to the cohort total).
The first subset comprises 37 examinations from 36 patients with whole-body
anatomical segmentation labels.
The second comprises 112 examinations from 95 patients with lesion segmentation
labels.
The third comprises 142 disease-free follow-up examinations from 95 patients,
each demonstrating complete metabolic response.
The anatomy-segmentation and lesion-segmentation subsets overlap in 29
examinations (32 patients at the patient level).
The anatomy-segmentation subset alone spans ages 1 to 17 years (median 11), and
is being expanded toward a larger target, so the released cohort is expected to
grow as further annotations are completed.
Inclusion criteria were [TBD: e.g., diagnostic-quality whole-body coverage], and
subjects were excluded for [TBD: reasons, with counts].
Participants represented a range of clinical presentations, both with and without
pathology, and the distribution is summarised in Figure~\ref{fig:cohort}.
A subject-by-subject description of demographic and clinical characteristics is
provided in the accompanying metadata (Section~\ref{sec:datarecords}).

\subsection*{Clinical MRI protocols}\label{subsec:acquisition}
% Reusable template: scanner(s)/coil/field strength, protocol variability by
% indication, shared sequence(s), parameter table, contrast vs non-contrast.
% NOTE our TODO: confirm which sequences are included (e.g., whether DWI is
% included, and whether 3D Slicer can filter by sequence) -- see CLAUDE.md.
MRI data were acquired using [TBD: scanner manufacturer/model(s)] operating at
[TBD: field strength(s)] at [TBD: site(s)]. Because imaging protocols were
determined by clinical indication, the number and type of sequences vary across
subjects, and include both contrast-enhanced and non-contrast acquisitions.
Whole-body coverage was obtained using [TBD: acquisition strategy, e.g.,
multi-station composed sequences], and acquired contrasts include [TBD: e.g.,
T1-weighted, T2-weighted, STIR, Dixon, DWI]. Acquisition parameters for the most
commonly used protocol are summarised in Table~\ref{tab:acq}; per-scan parameters
are provided in the accompanying sidecar JSON files
(Section~\ref{sec:datarecords}). [TBD: note any sequences deliberately excluded
and why --- cf. the reference, which excluded diffusion data owing to less
standardised quality assessment.]

%%-- Acquisition-parameter table placeholder (cf. Table 1 of ref) ----------%%
\begin{table}[h]
\caption{Representative whole-body MRI acquisition parameters for the most
frequently applied protocol (placeholder --- to be completed).}\label{tab:acq}
\begin{tabular}{@{}llll@{}}
\toprule
Parameter & [Seq.\ 1] & [Seq.\ 2] & [Seq.\ 3] \\
\midrule
TR [ms]            & [TBD] & [TBD] & [TBD] \\
TE [ms]            & [TBD] & [TBD] & [TBD] \\
Flip angle [$^\circ$] & [TBD] & [TBD] & [TBD] \\
Voxel size [mm]    & [TBD] & [TBD] & [TBD] \\
Orientation        & [TBD] & [TBD] & [TBD] \\
\botrule
\end{tabular}
\end{table}

\subsection*{Data processing}\label{subsec:curation}
% Highly reusable: DICOM->NIfTI conversion + organisation; de-identification.
% The reference used ezBIDS/dcm2niix + pydeface (an automatic de-facer) -- this
% is directly relevant to our de-identification TODO (CLAUDE.md). Skull-stripping
% (BET) is brain-specific and likely NOT applicable to whole-body; replace with
% our QC/masking approach as needed.
DICOM data were converted to NIfTI and organised following [TBD: BIDS-aligned /
custom open] conventions (Section~\ref{sec:datarecords}) using [TBD: tool, e.g.,
dcm2niix]. To de-identify the imaging, [TBD: describe the de-identification
pipeline]. Where head-inclusive coverage permits facial reconstruction,
identifiable facial features were removed using [TBD: an automatic de-facer,
e.g., a pydeface-style tool --- see the de-identification TODO in CLAUDE.md].
An overview of data collection, acquisition, and processing is presented in
Figure~\ref{fig:workflow}.

% Where radiology reports are included (Data Records), they are de-identified
% ("scrubbed") to remove protected health information before release; validation
% is reported in Technical Validation.
Where de-identified radiology reports are released with the imaging, protected
health information was removed using [TBD: report-scrubbing pipeline]; the
validation of this scrubbing is reported in Section~\ref{sec:validation}.

\subsection*{Annotation and labelling}\label{subsec:annotation}
% Reusable concept: expert review + structured labels. Our addition: labels keyed
% by reader and anatomy to enable inter-reader agreement (Technical Validation).
Labels are provided to support supervised model development and are organised so
that provenance is explicit: each label records the target anatomy, the label
type ([TBD: segmentation / bounding box / image-level classification]), and the
annotating reader together with their experience level. Annotation was performed
using [TBD: tool/version] following [TBD: protocol/atlas]. A subset of [TBD: $N$]
examinations was independently labelled by [TBD: $N$] readers to enable
inter-reader agreement analysis (Section~\ref{sec:validation}).

%%-- Workflow figure placeholder (cf. Fig. 2 of ref) -----------------------%%
\begin{figure}[t]
\centering
\fbox{\parbox[c][4.5cm][c]{0.85\textwidth}{\centering\itshape
[Figure placeholder: acquisition and processing workflow --- from imaging
through conversion, de-identification, labelling, and quality assessment.]}}
\caption{Overview of the data acquisition and processing workflow. [TBD:
finalise.]}
\label{fig:workflow}
\end{figure}
